\section{Recomendaciones para la nube}
Son criterios generales, no selección validada de proveedor o instancia.
\subsection{Dimensionamiento}
CPU y memoria deben contrastarse con throughput requerido, latencia,
estados y margen. Cuotas locales no equivalen automáticamente a vCPU
dedicadas, créditos o almacenamiento remoto.

RSS individual no define un mínimo seguro. No se recomienda reducir
memoria hasta el máximo observado ni aumentar CPU solo por throttling.
Deben contemplarse concurrencia, cachés y trabajo temporal.

\subsection{Concurrencia y escalabilidad}
Más clientes podrían aumentar colas sin mejorar QPS, pero no está
demostrado aquí. C10/C50 no reciben recomendaciones cuantitativas.

Escalar horizontalmente introduce coordinación, red y consistencia.
Los beneficios no se deducen de un contenedor local por motor.

\subsection{Selección del motor}
Q1 no representa toda la aplicación. Importan transacciones, escrituras,
modelo de datos, índices y operación. Dos observaciones no establecen
un ganador general. Q2 requiere evaluación y exposición de la limitación
del índice textual.

\subsection{Memoria y almacenamiento}
La memoria por operación debe revisarse con sesiones, límite del
contenedor y almacenamiento temporal. Aumentar parámetros puede trasladar
presión desde disco a memoria. Spill no implica automáticamente fallo;
su ausencia tampoco garantiza una ejecución óptima.

\subsection{Observabilidad}
Deben conservarse identificadores, commit, consultas, variantes, límites
y cobertura. Los paneles deben distinguir ausencia de datos y cero.
La reutilización del preflight requiere contexto compatible; su bloqueo
no permite afirmar ahorro demostrado.
